% SPDX-License-Identifier: MIT
\chapter{Installation}
\label{ch:installation}

\app{} runs on Windows, macOS and Linux from a single environment definition. The
commands are the same on all three; the environment resolves the platform
differences, including the Qt system libraries that would otherwise have to be
installed by hand on Linux.

\app{} is not on PyPI, so installation is from the source repository. There is no
\cmd{pip install makeyourtree} from the index --- ignore any package by that name
from another author.

\section{Requirements}

\begin{itemize}
  \item \cmd{git}
  \item A conda-family package manager, or Python 3.12 or newer for the pip route
  \item About 1\,GB of disk space. Most of that is Qt: a complete environment
        measures 751\,MB, of which PySide6 is 641\,MB and NumPy 34\,MB.
\end{itemize}

No network access is needed at runtime --- only to install.

\section{With conda or mamba}

This is the recommended route. Substitute \cmd{mamba} or \cmd{micromamba} for
\cmd{conda} throughout if you prefer; the arguments are identical and mamba
resolves considerably faster.

\begin{lstlisting}[language=bash]
git clone https://github.com/bioforkgithub/makeyourtree.git
cd makeyourtree
conda env create -f environment.yml
conda activate makeyourtree
pip install -e . --no-deps
\end{lstlisting}

\begin{quote}
\textbf{\opt{--no-deps} is not optional bookkeeping.} Conda has already installed
everything \app{} needs. A plain \cmd{pip install -e .} would resolve those same
requirements again from PyPI and install a \emph{second} copy of Qt into an
environment that already has one --- 641\,MB wasted, and two Qt builds present
where which one loads depends on import order.
\end{quote}

\href{https://github.com/conda-forge/miniforge}{Miniforge} is the usual
recommendation for the package manager itself: it is small and defaults to the
conda-forge channel, which is where \cmd{pyside6} comes from. On Anaconda or
Miniconda the \cmd{defaults} channel takes priority and does not carry
\cmd{pyside6}; if the environment fails to solve, prefer conda-forge:

\begin{lstlisting}[language=bash]
conda config --add channels conda-forge
conda config --set channel_priority strict
\end{lstlisting}

\section{With pip and venv}

Use this if you would rather not install conda. It needs Python 3.12 or newer
already present, and on Linux it needs the Qt system libraries installed by hand
(Section~\ref{sec:linux-libs}).

\begin{lstlisting}[language=bash]
git clone https://github.com/bioforkgithub/makeyourtree.git
cd makeyourtree
python -m venv .venv
\end{lstlisting}

Activate it. This is the only command in this manual that differs by platform:

\begin{center}
\begin{tabular}{@{}ll@{}}
\toprule
\textbf{Platform} & \textbf{Command} \\
\midrule
Linux, macOS         & \cmd{source .venv/bin/activate} \\
Windows, PowerShell  & \cmd{.venv\textbackslash Scripts\textbackslash Activate.ps1} \\
Windows, \cmd{cmd.exe} & \cmd{.venv\textbackslash Scripts\textbackslash activate.bat} \\
\bottomrule
\end{tabular}
\end{center}

Then install. Note that here pip \emph{does} resolve dependencies, so there is no
\opt{--no-deps}:

\begin{lstlisting}[language=bash]
pip install -e ".[studio]"
\end{lstlisting}

\section{The library alone, without Qt}
\label{sec:install-library}

The \cmd{makeyourtree} package and its command-line interface import no GUI
toolkit, so they install without Qt. This is the right choice for a server, a
container, a CI job or an analysis pipeline: it avoids the 641\,MB PySide6 brings
and needs none of the Linux system libraries below.

\begin{lstlisting}[language=bash]
pip install -e "."
\end{lstlisting}

You then have the \cmd{makeyourtree} command and the Python API, but not
\cmd{makeyourtree-studio}. NumPy is the only dependency. With conda, delete the
\cmd{pyside6} line from \cmd{environment.yml} before creating the environment.

\section{Checking the installation}

\begin{lstlisting}[language=bash]
makeyourtree --version
makeyourtree info examples/primates.nwk
\end{lstlisting}

To run the test suite as well --- the fastest way to confirm a healthy install:

\begin{lstlisting}[language=bash]
pip install -e ".[studio,dev]"
python -m pytest
\end{lstlisting}

Expect \textbf{2074 passed} in roughly 30 seconds. The suite runs headlessly and
needs no display.

\section{Platform notes}

Two things genuinely differ by operating system. Both affect the pip route only;
the conda environment handles the first and largely avoids the second.

\subsection{Linux: Qt system libraries}
\label{sec:linux-libs}

The PySide6 \emph{wheel} bundles Qt itself, but Qt still loads a few system
libraries --- above all its \cmd{xcb} platform plugin. Without them the
application exits with
\cmd{qt.qpa.plugin: Could not load the Qt platform plugin "xcb"}.

\begin{lstlisting}[language=bash]
# Debian / Ubuntu
sudo apt install libgl1 libegl1 libglib2.0-0 libfontconfig1 libdbus-1-3 \
  libxkbcommon-x11-0 libxcb-cursor0 libxcb-icccm4 libxcb-image0 \
  libxcb-keysyms1 libxcb-randr0 libxcb-render-util0 libxcb-shape0 \
  libxcb-xinerama0

# Fedora
sudo dnf install mesa-libGL mesa-libEGL fontconfig dbus-libs \
  libxkbcommon-x11 xcb-util-cursor xcb-util-image xcb-util-keysyms \
  xcb-util-renderutil xcb-util-wm
\end{lstlisting}

\cmd{libxcb-cursor0} is required by Qt~6.5 and later specifically, and is the one
most often missing on otherwise complete systems. On Debian and Ubuntu,
\cmd{python3.12-venv} is also a separate package.

\subsection{Windows: clone to a short path}

PySide6 ships files whose names exceed Windows' 260-character \cmd{MAX\_PATH}
limit. Installing into a \cmd{.venv} inside a deeply nested directory fails with
\cmd{OSError: [WinError 206] The filename or extension is too long}, and it rolls
back \emph{partially} --- which is confusing, because the test suite still passes
while \cmd{makeyourtree} itself is missing.

Clone somewhere short such as \cmd{C:\textbackslash dev\textbackslash makeyourtree},
or enable long paths with \cmd{git config -{}-system core.longpaths true} together
with the \cmd{LongPathsEnabled} registry setting. The conda route mostly sidesteps
this, because the environment lives under the conda installation rather than
inside the clone.

\section{Building a standalone bundle}

To produce a self-contained directory that runs without Python installed:

\begin{lstlisting}[language=bash]
pip install pyinstaller
python -m makeyourtree_studio.packaging.build
\end{lstlisting}

The result is \cmd{dist/MakeYourTreeStudio/}. The build driver verifies the
bundle against thirteen checks --- licence texts present, Qt shipped as separate
replaceable shared libraries, no GPL-3.0-only Qt add-on pulled in --- and refuses
to certify one that fails.

PyInstaller cannot cross-build: each platform's bundle must be produced on that
platform. At the time of writing this has been built and launched on Windows
only; macOS and Linux bundles are unproven rather than broken.
